\chapter{Introduction \& Clinical Context}

\section{Background \& The Clinical Crisis of Triage Delays}

Emergency departments (EDs) and outpatient triage facilities represent the high-stakes frontline of modern healthcare delivery. In these acute care environments, patient arrival patterns are inherently stochastic, exhibiting extreme volatility during infectious disease outbreaks, seasonal shifts, and community surges. Among all patient demographics, the pediatric population (spanning neonates, infants, toddlers, and young children) possesses the highest vulnerability to rapid physiological deterioration.

Unlike adult patients, who frequently communicate pain, localized symptoms, and progressive weakness verbally, pediatric patients exhibit non-specific clinical signs of distress. An infant experiencing respiratory syncytial virus (RSV), acute gastroenteritis with severe dehydration, or septicemia may transition from quiet lethargy to irreversible hemodynamic collapse within a narrow therapeutic window of 30 to 60 minutes. Consequently, timely identification, continuous triage queue surveillance, and prompt clinical escalation are paramount to preventing avoidable pediatric morbidity and mortality.

Despite the critical nature of early pediatric intervention, hospital waiting rooms in resource-constrained environments, particularly throughout developing healthcare infrastructures in Sub-Saharan Africa and underserved public municipal hospitals, suffer from chronic understaffing and manual patient tracking bottlenecks. Triage nurses are overwhelmed by manual paper-based logging, static intake registries, and intermittent visual scans of crowded waiting halls. Under such operational strain, hospital administrators lack real-time visibility into the exact number, proportion, and wait duration of pediatric patients in the waiting area.

\section{The "Invisible Child" Phenomenon: Definition \& Etiology}

A primary driver of pediatric triage oversight is a widespread yet under-researched clinical reality termed in this thesis as the \textbf{"Invisible Child" phenomenon}.

\begin{verbatim}
  +-------------------------------------------------------------------------+
  |                    THE INVISIBLE CHILD PHENOMENON                       |
  |                                                                         |
  |   1. Physical Occlusion:                                                |
  |      - Infant held against caregiver's chest (kangaroo care / sling)    |
  |      - Swaddled in thick blankets obscuring human contour               |
  |      - Blocked by waiting room chairs, luggage, or crowding             |
  |                                                                         |
  |   2. Acoustic & Behavioral Masking:                                     |
  |      - Sick infants become hypoactive / lethargic rather than crying    |
  |      - Quiet status mistaken by staff for peaceful resting              |
  |                                                                         |
  |   3. Triage Registry Blindspots:                                        |
  |      - Caregivers register under their own name without noting infants  |
  |      - Manual headcounts count 1 adult caregiver = 1 patient            |
  |                                                                         |
  |   ===> SYSTEMIC RESULT: Severe pediatric load undercounting             |
  |        Nurse-to-pediatric-patient ratios dangerously misallocated       |
  +-------------------------------------------------------------------------+
\end{verbatim}

The Invisible Child phenomenon arises from three compounding factors:

\begin{enumerate}
    \item \textbf{Severe Physical and Morphological Occlusion:}  
   In clinical waiting halls, infants and young children rarely sit upright in isolated chairs. Instead, they are carried against caregivers' torsos in fabric wraps or slings, held in nursing positions, seated low to the floor, or swaddled within thick blankets. From the vantage point of standard wall-mounted or ceiling surveillance cameras, the visual boundary between the adult caregiver and the carried child is heavily obscured, causing conventional object detection algorithms and human staff to register only a single adult entity.

    \item \textbf{Acoustic and Behavioral Masking:}  
   Critically decompensating pediatric patients often do not exhibit overt distress such as loud crying. In advanced stages of respiratory exhaustion, hypoxemia, or metabolic acidosis, pediatric patients become placid, somnolent, and motionless. To busy triage personnel scanning a room, a lethargic, dying infant appears identical to a sleeping, healthy child.

    \item \textbf{Intake Registry Disconnects:}  
   In busy public clinics, adult caregivers often register at intake counters under their own primary identification without immediate clerical registration of accompanying infants. Consequently, electronic medical records (EMR) or paper logbooks reflect adult attendance while completely missing the presence of critically ill pediatric dependents sitting in the waiting area.
\end{enumerate}

When pediatric patient volume is systematically undercounted, hospital administrators allocate nursing and clinical triage staff based on flawed occupancy statistics. If a waiting room containing 20 adults actually harbors 15 carried, uncounted infants, the pediatric nursing capacity is overwhelmed, leading to catastrophic delays in care.

\section{Problem Statement}

Modern healthcare facilities require automated, non-invasive, and real-time surveillance of waiting room occupancy to maintain safe nurse-to-patient staffing ratios and prevent pediatric neglect. However, deploying computer vision systems in resource-constrained clinical settings faces four formidable challenges:

\begin{enumerate}
    \item \textbf{Severe Occlusion, Scale Disparity \& Semantic Ambiguity:} Existing computer vision models trained on generic pedestrian benchmarks (such as COCO or Pascal VOC) fail to detect small, occluded, or non-standard pediatric postures, generating high false-negative rates ($>70\%$) for carried infants. Standard convolution heads lack the dense semantic context required to distinguish folded blankets from swaddled infants.
    \item \textbf{Identity Fragmentation \& Counting Duplication:} In crowded, dynamic waiting halls, patients move, stand, sit, and temporarily pass behind barriers. Classical Multi-Object Tracking (MOT) algorithms frequently lose track continuity during prolonged occlusion, assigning new IDs upon target reappearance and causing severe cumulative double-counting in daily patient statistics.
    \item \textbf{High Latency \& Hardware Cost:} High-parameter Vision Transformers (ViTs) and large foundation models require power-hungry discrete GPUs ($>\$1,500$) and complex cloud streaming backends. Resource-constrained rural and municipal clinics cannot support expensive hardware infrastructure, necessitating lightweight edge-executable architectures that operate on standard consumer-grade CPUs.
    \item \textbf{Network Serialization \& Web Lag:} Conventional distributed architectures (separating computer vision microservices from web interfaces via REST polling or heavy WebSocket frameworks) introduce buffer backpressure, frame drops, and latency, resulting in frozen feeds and desynchronized telemetry.
\end{enumerate}

There is an urgent necessity for a \textbf{computationally lightweight, illumination-invariant, occlusion-robust, and edge-deployable pediatric monitoring system} that leverages self-supervised foundation knowledge distillation, robust tracking continuity, and instant clinical telemetry on low-cost hardware.

\section{Research Questions \& Core Hypotheses}

\subsection{Research Questions}
This thesis investigates the following fundamental research questions:

\begin{itemize}
    \item \textbf{RQ1:} Can dense feature-based knowledge distillation from self-supervised Vision Foundation Models (\textbf{DINOv3 ViTs}) empower lightweight edge detectors (\textbf{YOLO26s}) to resolve swaddled and heavily occluded pediatric patients without edge CPU inference penalties?
    \item \textbf{RQ2:} To what extent does Slicing Aided Hyper Inference (\textbf{SAHI}) resolve small-scale infant bounding boxes in wide-field, high-resolution hospital CCTV camera views?
    \item \textbf{RQ3:} How can dynamic multi-tracking algorithms (ByteTrack, BoT-SORT, OC-SORT, FastTracker) be adaptively orchestrated alongside spatial centroid debouncing to eliminate ID fragmentation and prevent double-counting in crowded clinical triage environments?
    \item \textbf{RQ4:} To what extent does LAB color space Contrast Limited Adaptive Histogram Equalization (\textbf{CLAHE}) improve feature saliency and detection recall in poorly illuminated clinical waiting rooms without inducing chromatic distortion?
    \item \textbf{RQ5:} Can an optimized, quantized ONNX Runtime pipeline running on a consumer-grade CPU match or exceed the operational throughput (30 FPS) of heavy GPU-dependent architectures while sustaining real-time reactive UI telemetry?
\end{itemize}

\subsection{Core Hypotheses}
\begin{quote}
\textit{It is hypothesized that transferring rich spatial representation priors from a DINOv3 Vision Transformer teacher to a YOLO26s convolutional student via cosine feature alignment and MSE loss, coupled with SAHI patch slicing, LAB CLAHE normalization, and an adaptive multi-tracker suite within a zero-network-serialization monolithic architecture, will achieve $>94\%$ pediatric detection precision, $>75\%$ recall under severe occlusion ($>3\times$ baseline COCO recall), $<5\%$ ID fragmentation, and real-time ($>30\text{ FPS}$) edge CPU execution, thereby providing a reliable automated safeguard against the Invisible Child phenomenon in clinical triage environments.}
\end{quote}

\section{Research Objectives}

To validate the hypothesis and address the research questions, the project defines the following specific objectives:

\subsection{Primary Objective}
To design, implement, mathematically formulate, and evaluate an edge-deployed computer vision and dynamic multi-tracking monitoring system that provides real-time pediatric capacity intelligence and overcrowding alerts in clinical triage waiting areas.

\subsection{Specific Objectives}
\begin{enumerate}
    \item \textbf{Foundation Knowledge Distillation \& Model Optimization:} Design a 2-stage training regimen transferring semantic patch representations from a DINOv3 ViT teacher to a YOLO26s student using a $1\times 1$ conv feature projection head and joint cosine/MSE loss, exporting optimized ONNX runtime weights for CPU edge acceleration.
    \item \textbf{Slicing Aided Hyper Inference (SAHI) Integration:} Formulate a multi-scale patch tiling and Non-Maximum Suppression (NMS) reconstruction engine to resolve small-scale carried infants in wide-angle triage CCTV feeds.
    \item \textbf{Illumination Normalization Pipeline:} Implement a real-time LAB color-space CLAHE preprocessing algorithm to enhance structural contrast in dim triage environments without perturbing color channels.
    \item \textbf{Adaptive Multi-Tracker Suite \& Scene Analyzer:} Develop an intelligent tracking engine that dynamically evaluates frame-to-frame optical flow motion ($\Delta I$) and crowd occlusion density ($\text{IoU}_{\text{pairwise}}$), automatically routing frames between ByteTrack, BoT-SORT, OC-SORT, and FastTracker with a 3.0-second hysteresis stabilization barrier.
    \item \textbf{Spatial Fallback \& Temporal Debouncing:} Formulate an untracked spatial centroid re-identification mechanism ($r \le 40\text{ px}$) and a 5.0-second lost-track debouncing queue to suppress identity switching and prevent cumulative patient overcounting.
    \item \textbf{Reactive Monolithic Clinical Dashboard:} Build a high-performance, single-process Python application using NiceGUI that binds computer vision inference state directly to clinical UI components, providing live stream visualization, interactive model/camera switching, capacity threshold alerting, and automated CSV/PDF clinical audit reporting.
    \item \textbf{Empirical Benchmarking \& 6-Way Ablation Evaluation:} Conduct rigorous empirical evaluations across a 6-way comparative ablation matrix ($M_1 \dots M_6$), measuring COCO 101-point mAP@50, mAP@[50:95], occlusion-tier AP, distillation fidelity, MOTA, IDF1, latency percentiles (p50/p95), and CPU/GPU utilization.
\end{enumerate}

\section{Scientific \& Engineering Contributions}

The primary contributions of this thesis are summarized as follows:

\begin{itemize}
    \item \textbf{Theoretical \& Algorithmic Contribution (Foundation Representation Distillation):} Formulated and verified a dense multi-scale feature distillation loss combining Cosine Similarity and Mean Squared Error ($\mathcal{L}_{\text{distill}} = \alpha \mathcal{L}_{\text{cos}} + \beta \mathcal{L}_{\text{MSE}}$) between a DINOv3 ViT teacher and YOLO26s student, boosting severe occlusion recall from \textbf{24.8\% (COCO Base)} and \textbf{48.2\% (Fine-Tuned)} to \textbf{76.9\% (Distilled + SAHI)}.
    \item \textbf{Computer Vision Pipeline Contribution (SAHI \& CLAHE Fusion):} Engineered a combined multi-scale patch slicing (SAHI) and LAB color-space CLAHE pipeline that enhances small-scale infant saliency in low-light hospital corridors without chromatic aberration.
    \item \textbf{Tracking Contribution (Adaptive Multi-Tracking \& Parent-Child Anchoring):} Introduced an adaptive tracker selection engine governed by live scene analysis metrics (mean frame difference optical flow and pairwise IoU overlap) paired with a Euclidean spatial fallback mechanism, reducing tracking ID switches by \textbf{68.5\%} compared to baseline SORT tracking.
    \item \textbf{Architectural Contribution (Zero-Latency Decoupled Monolith):} Formulated an asynchronous dual-thread execution model (\texttt{CAP\_PROP\_BUFFERSIZE = 1}) integrated into a NiceGUI in-memory reactive state architecture, eliminating inter-process network serialization overhead and ensuring zero-delay 30 FPS video streaming on consumer hardware.
    \item \textbf{Performance Contribution (ONNX Edge Quantization):} Achieved a \textbf{2.1x CPU inference acceleration} (28.4 ms/frame vs. 59.8 ms/frame on PyTorch FP32), establishing the feasibility of running full clinical AI surveillance on modest edge hardware without dedicated GPUs.
    \item \textbf{Clinical \& Empirical Contribution (Governance \& 6-Way Matrix):} Delivered an automated clinical capacity governance framework ($C_{\text{child}}\%$) with automated alert dispatching and generated publication-ready LaTeX \texttt{booktabs} ablation tables evaluating all 6 system configurations.
\end{itemize}

\section{Thesis Organization}

The remainder of this thesis is structured as follows:

\begin{itemize}
    \item \textbf{Chapter 2 (Literature Review):} Reviews the mathematical and historical evolution of deep learning object detectors, foundation models, knowledge distillation paradigms (Hinton to DINOv3), SAHI patch inference, multi-object tracking (SORT to ByteTrack), histogram equalization, and perspective geometry.
    \item \textbf{Chapter 3 (Methodology \& System Architecture):} Details the monolithic NiceGUI system design, offline distillation pipeline, SAHI architecture, and the mathematics of the adaptive Multi-Tracker Suite.
    \item \textbf{Chapter 4 (Implementation, Experiments \& Results):} Presents the clinical telemetry monitoring UI, followed by the experimental methodology and a 6-way comparative ablation matrix evaluating detection and tracking performance across all system configurations.
    \item \textbf{Chapter 5 (Conclusion \& Future Work):} Discusses ethical deployment, privacy-by-design, summarizes key findings, highlights clinical impacts, and charts future research pathways.
\end{itemize}